\chapter{Parallels}
``How was school today?'' Shreya asked.

``It was pretty good. I found out I got the highest marks in Maths,'' Shiva replied, holding the phone closer, tightening his grip on the pillow he had already squeezed into a ball.

``That's great,'' Shreya said. ``I wish I could do that well in Maths. I'm so worried about my boards. They say they'll be in February, and it's November already. I swear, if I just scrape a pass in this subject, I'll never touch it again. I'd rather take commerce in eleventh.''

``Why do you worry so much? It's easy once you get the core,'' Shiva said, half joking.

``You don't know it. Next year, when your boards come, I'll ask you then. I have pre-boards next month. If I get low marks, I don't know what Mummy will do to me.'' Shreya's voice grew audibly anxious.

``Don't worry. It'll be all right.''

``I don't know how you can be so relaxed about all this. For us girls, there are two options—marry or study—and I'm not interested in the first one yet.''

``Why, what's so wrong with that? Just imagine—you wouldn't have to worry about earning, working outside, dealing with the world. Isn't that a good option?''

``You can say that because you have a choice—you're a boy. You can't understand what girls have to go through in this country. If you fail your exams, you'll get married; if your family isn't wealthy enough, you'll get married; and if you happen to love someone, you'll get married too—to someone else. I wish I'd been born somewhere else, not here.''

``Okay, my princess. Don't overstress yourself. It's just a board exam, and no one is going to marry you off to someone else. Besides, I trust you'll do really well.'' Shiva smiled and tugged the pillow tighter, as if it might disappear any moment. ``But hey, tell me one thing—is there any plan of you visiting me soon?''

``I think so. I heard Bhaiya will be coming to your house next month for the Pooja. I think I'll come along too.''

``What? There's a pooja in my own house and I don't even know about it,'' Shiva exclaimed.

Shreya laughed. ``Apparently, kids don't get to know everything.''

``I'm not a kid. I'm just a year younger than you!'' Shiva shot back.

``Yeah, coming!'' Shreya called out suddenly—someone must have been calling for her.

``What happened? Who's calling you now?'' Shiva's voice dropped, a little crestfallen.

``Boy, I've got some chores to do.''

``Wait—what chores at three in the afternoon?''

``I have to mop the hallway. Dadi asked me this morning, so now that I'm back from school I have to finish it before evening—after that I have to join Mummy for cooking dinner.''

``Oh god. Can't you stay a little longer? It hasn't even been half an hour, and you just got back from school. Get some rest,'' Shiva said, and then his voice turned playful. ``And some time with me.''

``No, sorry. I really have to—otherwise, if Mummy comes upstairs and sees me on a call with someone, it'll turn into my marriage day,'' she laughed.

``Okay, fine. Bye—take care. And I love you,'' he said with a sigh.

``Bye-bye. Love you more.''

He set the phone aside, finally letting go of the pillow, and stretched out on the bed, looking up. The fan was still now, and it was a little cold outside. Light from the exhaust vent crept into the dark room and fell across the motionless blades, lighting up the dust that floated in the air. Shiva blew hard, and the dust scattered into a mesmerising pattern. The light felt as though it were spilling from some heaven's door in a television serial—as if a fairy might step through it at any moment.

Late in the evening, Shiva returned from being out with his friends.

``Get ready, freshen up. Dinner's ready,'' Nani called out the moment he stepped in through the gate.

``Yes, Nani, just a few minutes.''

He came back with a bottle of hair oil and settled beside Nani while Nana ate his dinner.

``You haven't asked for the coaching fees. It's already the second of November,'' Nana said.

``Yes, I was thinking about it. I found out that one of the teachers from my school has started teaching after hours. I asked him if I could join there instead, and he agreed. I think it'll be much better—now I won't have to go five kilometres out in the cold mornings.''

``How is he? You told me Robin was a good teacher, that a lot of students go to him.''

``He's good too. In fact, he lives in the school itself—the school hired him from Bengal, and he's really good at Mathematics and Science,'' Shiva replied, slowly handing the oil bottle to Nani and settling in front of her, a little lower.

Nani quietly worked the oil into her palm and began massaging it into his head, as if she already knew just what to do.

``Fine. As long as you keep working hard, the place doesn't matter. It's you who has to study, not the teacher,'' she added.

``Yes. Sure, Nani,'' Shiva nodded.

``How much is his fee?'' Nana asked.

``It's even less than Robin sir's. He only takes five hundred—two hundred less—and teaches two subjects instead of one.''

``Okay. But go to Robin tomorrow morning and pay last month's fees. And when are you joining the teacher at school?''

``From tomorrow,'' Shiva replied, slowly sinking down and resting his head in Nani's lap.

``Mummy ji, shall I bring your dinner too?'' Mami asked from the gallery nearby.

``Sure,'' Nani said. ``And Shiva, get me some water to wash this oil off my hands first.''

Shiva got up slowly. ``Sure, Nani—just a minute.''

The next day, as soon as school let out, Shiva walked straight up to the second floor, his bag still on his shoulders. Roshini was standing by the railing, just outside the door of the room they were to study in.

``Hello, Roshini. How are you?'' Shiva asked, all enthusiasm.

``Oh, hello. I'm great. What are you doing here?'' she asked, surprised.

``Were you looking for me on the court?'' he replied with a smile.

``No, no... Why would I look for you? But seriously, what are you doing here?'' she asked again.

``Just checking whether you're doing your best in your studies,'' he said, his tone turning teasing.

``Don't tell me you've joined the tuition with us?'' Her mouth fell open.

``Sure did,'' he said, lowering his head with a smile. ``Now that Mr. Verma himself is teaching, how can I not attend?''

He peeked through the open door into the room.

``She's inside. Just having her lunch,'' Roshini added, teasing.

``No... no... I wasn't looking for her. I was only checking if sir had arrived.''

``Do you think I'd be standing out here if sir were already in?''

``That's true too. Shall we go in?'' he asked.

They both walked in. The front seat was turned around to face the rest—Mr. Verma's, most likely—and in the seat before it sat Gauri, her lunchbox open.

Shiva crossed to her, threw his arms open in a film-hero flourish, and grinned. ``Hi—I'm Shiva. You must have heard the name.'' He held out his hand for a handshake.

Everything went still. No one moved, and Shiva stood there frozen, hand out, waiting. Then, after a moment, Gauri simply lifted her own hand, fingers glistening mid-bite, to show she was eating.

``Oh. Right. My bad.'' He pulled the hand back and, to save face, shook it firmly with his own.

Roshini burst out laughing. ``Too bad you're not Rahul—the one Shahrukh plays. That entrance only works in the movies.''